\documentclass[11pt]{article}

\usepackage[letterpaper,margin=1in]{geometry}
\usepackage{amsmath,amssymb}
\usepackage{booktabs}
\usepackage{array}
\usepackage{enumitem}
\usepackage{xcolor}
\usepackage[hidelinks]{hyperref}

\setlength{\parindent}{1.5em}
\setlength{\parskip}{0pt}
\renewcommand{\arraystretch}{1.15}

\begin{document}

\begin{center}
    {\LARGE Project 1: Biomedical Signal Processing}\par
    \vspace{1.5em}
    {\large Due: Sunday, 1st November, 2026, 23:59}\par
    {\large Responsible TA: Yige Gao}\par
    {\large Email: \href{mailto:gaoyg2025@shanghaitech.edu.cn}{gaoyg2025@shanghaitech.edu.cn}}
\end{center}

\vspace{1.5em}

\section{Introduction}

The evolving modern lifestyle has led to an increasing prevalence of lumbar spine disorders. Epidemiological studies indicate that up to 80\% of the population experiences spinal ailments at various stages of their lifetime. Accurate tissue identification during spinal procedures is crucial for surgical success and positive patient outcomes. Ultrasound signal analysis provides a fast and accurate method for tissue characterization, offering significant value for clinical applications.

This project investigates computational methods for classifying spinal tissues based on ultrasound signatures, integrating signal processing techniques with machine learning classification algorithms.

\section{Project Tasks}

\subsection{Dataset Description}

We will provide ultrasound signals collected from \textbf{four types of tissues} from sheep spines, together with their corresponding labels. The dataset is divided into \textbf{training, testing, and validation} sets.

All 14 signals belonging to the same biological sample are kept together within each data subset. This grouping must also be preserved in all subsequent training, validation, and evaluation procedures to prevent data leakage.

Each biological sample consists of 14 ultrasound signals, which are generated by a 16-element transducer. In the signal acquisition process, every 3 consecutive transducer elements are used together to form one ultrasound beam, resulting in a total of $16-3+1=14$ ultrasound beams per acquisition.

Each ultrasound signal contains \textbf{896 sampling points}. All signals are sampled at a frequency of $f_s=75\,\text{MHz}$. The raw signal amplitudes are stored in the range $[0,255]$.

Each of the 14 signals in one sample is associated with one corresponding tissue label. The four tissue categories and their label assignments are listed as follows:

\begin{itemize}
    \item Spinal cord: 0
    \item Nucleus pulposus: 1
    \item Adipose tissue: 2
    \item Nerve root: 5
\end{itemize}

The detailed organization, file names, data formats, and dimensions of the training, test, and validation sets are summarized in Table~\ref{tab:dataset}.

\begin{table}[htbp]
\centering
\caption{Description of Dataset Files}
\label{tab:dataset}
\small
\begin{tabular}{@{}p{0.20\textwidth}p{0.12\textwidth}p{0.19\textwidth}p{0.41\textwidth}@{}}
\toprule
\textbf{File} & \textbf{Content} & \textbf{Size} & \textbf{Details} \\
\midrule
\texttt{train\_data.npy} & signal & $[506\times14,896]$ & Each row represents one signal. The length of each signal is 896, and there are a total of 7,084 signals. \\
\addlinespace
\texttt{train\_label.npy} & label & $[506\times14]$ & Each row corresponds to the label of the signal. \\
\addlinespace
\texttt{test\_data.npy} & signal & $[126\times14,896]$ & Each row represents one signal. Every 14 rows represent signals from the same sample. The length of each signal is 896, and there are a total of 126 samples. \\
\addlinespace
\texttt{test\_label.npy} & label & $[126\times14]$ & Every 14 rows correspond to the labels from the same sample. \\
\addlinespace
\texttt{val\_data.npy} & signal & $[126\times14,896]$ & Each row represents one signal. Every 14 rows represent signals from the same sample. The length of each signal is 896, and there are a total of 126 samples. \\
\addlinespace
\texttt{val\_label.npy} & label & $[126\times14]$ & Every 14 rows correspond to the labels from the same sample. \\
\bottomrule
\end{tabular}
\end{table}

\subsection{Preprocessing}

Each raw signal has amplitude values in the range $[0,255]$ and contains a DC offset of 127.5. Since the acquisition hardware cannot record decimal values, samples corresponding to the DC level may appear as either 127 or 128.

Before normalization, all values equal to 127 or 128 are replaced with 127.5. Then, each signal is normalized to the range $[-1,1]$ using:
\begin{equation}
    x_{\mathrm{norm}}[n]=\frac{x_{\mathrm{us}}[n]-127.5}{127.5},
\end{equation}
where $x_{\mathrm{us}}[n]$ denotes the original ultrasound signal and $x_{\mathrm{norm}}[n]$ is the normalized signal.

\subsection{Task 1: Feature Extraction and Statistical Classification}

Participants should employ signal processing methods to extract discriminative features from raw ultrasound signals. The \textbf{Empirical Wavelet Transform (EWT)}~\cite{gilles2013ewt} is recommended for its effectiveness in adaptive signal decomposition. Other traditional signal analysis and feature extraction methods, such as the \textbf{Short-Time Fourier Transform (STFT)}, \textbf{Discrete Wavelet Transform (DWT)}, and \textbf{Variational Mode Decomposition (VMD)}, may also be used, provided that the choice of method and parameter settings is justified.

After signal transformation or decomposition, students should extract statistical features from the resulting signal components or representations. When using EWT, decompose each signal into \textbf{multiple modes} and explain the number of modes and the rationale for its selection in the report. For each mode, extract at least the \textbf{mean, variance, median, kurtosis, root-mean-square (RMS), and skewness}, and combine the features from all modes into a feature vector for each signal. Additional statistical features may be included if their selection is justified.

\textbf{The same principle applies to other signal analysis and feature extraction methods:} For STFT, statistical features may be extracted from frequency bands or time-frequency regions of the magnitude or power representation. For DWT, features may be extracted from approximation and detail coefficients at the selected scales, or from the corresponding reconstructed components. For VMD, features may be extracted from the decomposed modes. Organize the resulting features into feature vectors for classification. The report must describe the signal representation, key parameters (e.g., the window function, window length, and overlap ratio for STFT; the wavelet and decomposition level for DWT; or the number of modes and penalty parameter for VMD), how components or regions are defined, and the quantities over which each statistical feature is computed, together with its definition.

Using the training set, students should compare the feature distributions across the four tissue types, identify discriminative features or feature combinations, and construct an explicit and interpretable \textbf{statistical classification rule}. The rule may use feature thresholds, value ranges, statistical scores, or other criteria based on statistical features, but it must clearly define how input features are mapped to tissue classes. Use the training set to estimate thresholds and statistical parameters; the validation set may be used to select or refine the classification rule. The test set may be used only for final evaluation after the method has been finalized, and must not be used to develop or adjust the classification criteria.

The report must describe the selected statistical features, the observed differences between tissue types, and the final classification rule and its parameters. Report classification results according to the sample-level evaluation strategy specified in this project. This task should use interpretable statistical criteria rather than directly applying the general-purpose machine learning classifiers used in Task 2.

\subsection{Task 2: Machine Learning Classification}

Using the features extracted in Task 1, students are required to implement machine learning classifiers for tissue classification.

Key requirements include:

\begin{itemize}
    \item Select appropriate ML models (e.g., SVM~\cite{cortes1995svm}, Random Forest~\cite{breiman2001rf}, k-NN~\cite{cover1967nn}).
    \item Perform \textbf{five-fold cross-validation} on the training set for model selection and hyperparameter tuning. The folds must be created at the \textbf{biological-sample level}: all 14 signals from one sample must remain in the same fold. Where feasible, use stratified folds to maintain a similar class distribution across folds.
    \item Conduct the complete feature-processing and model-fitting pipeline independently within each training fold. Any data-dependent preprocessing, feature selection, or normalization must be fitted using only the training portion of that fold and then applied to its held-out portion.
    \item Report the cross-validation performance as the \textbf{mean and standard deviation across the five folds}. Clearly state the fold-construction method and all hyperparameters considered.
    \item After selecting the model and hyperparameters through cross-validation, retrain the final model on the complete training set and evaluate it on the provided validation and test sets. Neither the validation set nor the test set may be included in the five-fold cross-validation or used to fit any preprocessing step.
    \item \textbf{Test-set restriction:} The test set must remain an independent holdout set. Neither the test data nor the test labels may be used for feature design, preprocessing, model selection, hyperparameter tuning, classification-threshold selection, early stopping, or any other form of model optimization. The test set may be used only after the model, hyperparameters, and complete processing pipeline have been finalized and training has been completed, solely to evaluate and report the model's final generalization performance.
    \item \textbf{Crucial:} Report classification results at the \textbf{sample level} (14 signals per sample), not at the individual signal level.
\end{itemize}

\subsection{Task 3: Deep Learning Classification (Vision Transformer)}

For the deep learning task, the \textbf{Vision Transformer (ViT)} model will be adopted. A trained network model will be provided. \textbf{Completing inference with the provided ViT model is mandatory}: every student must use this model for inference, report the classification results according to the sample-level evaluation strategy specified in this project, and compare its performance with the statistical and machine learning methods. Completing the bonus task described below does not replace this mandatory task.

\textbf{Independent-training bonus (10 points):} Students may independently train or fine-tune a ViT model using the provided data for an additional 10 points. The provided training, validation, and test split must be followed: the training set is used for model training, the validation set may be used for model selection, hyperparameter tuning, and early stopping, and the test set may be used only for final performance evaluation after all training and model-selection decisions have been completed. The existing dataset is sufficient for this task. The submission must include reproducible training code, the training configuration, a clear description of how each data subset was used, evaluation results, and a comparison with the provided model. Neither the test data nor the test labels may be used for training, model selection, hyperparameter tuning, or any other form of model optimization.

This bonus is awarded in addition to the original 100-point base score and does not change the base grading rules.

\subsection{Evaluation Metrics}

Classification performance should be evaluated using a confusion matrix (Table~\ref{tab:confusion}) and derived metrics.

\begin{table}[htbp]
\centering
\caption{Confusion Matrix for Binary Classification}
\label{tab:confusion}
\begin{tabular}{lll}
\toprule
& \multicolumn{2}{c}{\textbf{Predicted Class}} \\
\cmidrule(lr){2-3}
\textbf{Actual Class} & \textbf{Positive (P)} & \textbf{Negative (N)} \\
\midrule
\textbf{Positive (P)} & True Positive (TP) & False Negative (FN) \\
\textbf{Negative (N)} & False Positive (FP) & True Negative (TN) \\
\bottomrule
\end{tabular}
\end{table}

The following metrics should be calculated and discussed in your analysis:

\begin{enumerate}[label=\roman*).]
    \item \textbf{Accuracy:} Overall classification correctness
    \begin{equation}
        \mathrm{Accuracy}=\frac{TP+TN}{TP+TN+FP+FN}
    \end{equation}

    \item \textbf{Sensitivity (Recall):} Ability to identify positive instances
    \begin{equation}
        \mathrm{Sensitivity}=\frac{TP}{TP+FN}
    \end{equation}

    \item \textbf{Specificity:} Ability to identify negative instances
    \begin{equation}
        \mathrm{Specificity}=\frac{TN}{TN+FP}
    \end{equation}
\end{enumerate}

It should be noted that the above definitions of Accuracy, Sensitivity (Recall), and Specificity are based on a \textbf{binary classification setting}. In the binary case, the confusion matrix contains four basic elements: true positive (TP), false positive (FP), false negative (FN), and true negative (TN).

In this project, however, the tissue classification problem is formulated as a \textbf{four-class classification task}. To extend the evaluation metrics from the binary case to the multi-class case, a \textbf{one-versus-rest} strategy is adopted. Specifically, for each class, the corresponding $TP_c$, $FP_c$, $FN_c$, and $TN_c$ are computed by treating the current class as the positive class and all other classes as the negative class.

After obtaining the confusion statistics for each class, the overall performance of the four-class classification is evaluated using the \textbf{micro-averaging} strategy, as described below.
\begin{equation}
\begin{aligned}
    TP_{\mathrm{micro}}&=\sum_{c=1}^{C}TP_c, &
    FP_{\mathrm{micro}}&=\sum_{c=1}^{C}FP_c, \\
    FN_{\mathrm{micro}}&=\sum_{c=1}^{C}FN_c, &
    TN_{\mathrm{micro}}&=\sum_{c=1}^{C}TN_c.
\end{aligned}
\end{equation}
The overall evaluation metrics are then computed based on these aggregated values.

\subsection{Sample-level Evaluation Strategy}

Each biological sample contains 14 ultrasound signals. Therefore, classification results must be reported at the \textbf{sample level} rather than at the individual signal level. The mapping from 14 signal-level predictions to one sample-level label is performed using the following strategies.

\subsubsection{Voting Strategy (Recommended)}

Each signal is first classified independently, producing 14 predicted labels:
\begin{equation}
    \{y_1,y_2,\ldots,y_{14}\}.
\end{equation}
The final sample-level label $Y$ is determined by majority voting:
\begin{equation}
    Y=\arg\max_c\sum_{i=1}^{14}\mathbf{1}(y_i=c),
\end{equation}
where $\mathbf{1}(\cdot)$ is the indicator function.

\subsubsection{Averaging Strategy}

Let $\mathbf{f}_i$ denote the feature vector extracted from the $i$-th signal. The averaged feature vector of one sample is computed as:
\begin{equation}
    \bar{\mathbf{f}}=\frac{1}{14}\sum_{i=1}^{14}\mathbf{f}_i.
\end{equation}
The averaged feature vector is then used for sample-level classification.

Students are encouraged to explore more advanced evaluation strategies. All adopted methods must be clearly described in the final report.

\subsection{Discussion and Interpretation}

Critically analyze your classification results considering:
\begin{enumerate}[label=\roman*).]
    \item Performance differences between tissue types
    \item Effectiveness of your feature extraction method
    \item Limitations of your approach and potential improvements
\end{enumerate}

\section{Submission Requirements}

Submit a single zip file containing:
\begin{enumerate}
    \item \textbf{PDF Report:} Comprehensive document describing your methodology, results, and analysis. LaTeX source code or Word will \textcolor{red}{not} be accepted.
    \item \textbf{Source Code:} Your code with a README.md explaining how to run it.
    \item \textbf{Classification Results:} Your results should include the confusion matrix and the corresponding metrics.
\end{enumerate}

\textbf{Important:} Missing any component will result in at least 20\% grade deduction. Ensure your code is reproducible and well-documented.

\section{Grading Rules}

\begin{enumerate}
    \item Running the code successfully. (30pts)
    \item Accuracy of your results. (20pts)
    \item Report. (50pts)
\end{enumerate}

\noindent\textbf{Bonus:} Independently train or fine-tune a ViT model and satisfy the submission and evaluation requirements listed in Task 3. (10pts)

\begin{thebibliography}{9}

\bibitem{gilles2013ewt}
J\'er\^ome Gilles.
Empirical wavelet transform.
\emph{IEEE Transactions on Signal Processing}, 61(16):3999--4010, 2013.

\bibitem{cortes1995svm}
Corinna Cortes and Vladimir Vapnik.
Support-vector networks.
\emph{Machine Learning}, 20(3):273--297, 1995.

\bibitem{breiman2001rf}
Leo Breiman.
Random forests.
\emph{Machine Learning}, 45(1):5--32, 2001.

\bibitem{cover1967nn}
Thomas Cover and Peter Hart.
Nearest neighbor pattern classification.
\emph{IEEE Transactions on Information Theory}, 13(1):21--27, 1967.

\end{thebibliography}

\end{document}
